\chapter[Sermon 63. The Angel Preaches the Eternal Gospel of Christ.]{The Angel Preaches the Eternal Gospel of Christ.}
\label{sermon:063}
\markright{Sermon 63. The Angel Preaches the Eternal Gospel of Christ.}

And I saw an angel flying in the midst of heaven, having an everlasting gospel to preach to those who sit and dwell on the earth, and to every nation, kindred, language, and people, saying with a loud voice: Fear God, and give honor to him, for the hour of his judgment has come. Worship him who made heaven and earth, and the sea and the fountains of water.

\index{Antichrist}Antichrist desires nothing so much as to be oppressed by the preaching of the Gospel. For even therefore has he instituted inquisitors of heretical depravity to dare call the Gospel heresy. Therefore he burns Gospel books and preachers of the Gospel, and everywhere restrains the reading of the Gospel and Evangelical books. Wherefore the simple suppose that it can not be but that the Gospel with all its adherents should perish utterly. Now therefore, in the Lord’s consolation, is brought in a vision of an Angel (for he is still in the vision), flying in the midst of heaven, having the everlasting gospel, and preaching to the world. Whereby is signified that the Gospel shall be preached unto men, in spite of all the enemies thereof. And he gathers a brief summary of such things as by the Gospel are preached to the world. Those same appertain also to the comfort of the church, which under the old beast suffered persecutions for the Gospel. We will briefly consider every thing.

First, it is evident even from what has been said that “angel” signifies the ministers of the word, and the ministry of the gospel itself. Scripture calls preachers “angels.” John the Baptist is so named from the prophet Malachi, as was explained earlier. By this honorable title, ministers are reminded of purity and most sincere faith. For angels are God’s ministers, who regard, love, and honor only Him, and execute His commands most faithfully, sincerely, and diligently. Preachers should be like this in their kind and office. Just as angels cannot be harmed by the treasons and injuries of men, so God defends His ministers until the appointed hour. Peter was delivered from prison in Acts 12. Paul was protected in the shipwreck, and so on. And another angel is mentioned, because he has already introduced various visions of different angels. Nevertheless, “another” seems to refer to the first, for he adds two more angels, the first of whom he calls “another,” and the later, “the third.”

And this angel flies in the midst of heaven. By this is signified the fortunate course and progression of the preaching of the gospel. It is also written in the prophets, that his word runs swiftly. Psalm 19. David compares the running of the gospel’s preaching to the course of the sun, joyful as a giant he runs his way: he arises to the farthest part of the heavens and runs again to the same, and neither can any man stop him, nor hide himself from his heat. The sun shines in all places. Therefore, the preaching shall be free. For as we cannot pluck back or hinder the things that are above us in air and sky, so shall we neither pluck down nor hinder him that flies in the midst of heaven. The words and writings fly, they fly far and wide. Neither can the truth be oppressed. God has given to the world Printing, whereby the gospel is preached and runs far, wide, and most swiftly.

And this angel has the everlasting gospel. Wherein is the greatest comfort. For it signifies that the truth will be in the world invincible. And for many reasons the gospel is called everlasting. First, because the truth is immortal, which cannot be bound, however the ministers are fettered and slain. 2 Timothy 2. Secondly, the gospel is eternal, for because it was shown to our first fathers, prophesied in the law and prophets, fulfilled by Christ, declared by the apostles, and by the grace of God brought unto us. Indeed, it was predestined before all times. Read 1 Ephesians. For this cause it is called everlasting, as much as it pertains to us and to our posterity until the end of the world, and not only to our elders. And because it is everlasting, those lie who at this day call it a new doctrine or learning. Papistry is new, which has its origin when everything was ordained. Moreover, the Apostle says: If I or an angel from heaven should preach any other gospel, or besides the one that you have received, let him be accursed.

\index{Daniel (Book of)}And we hear expressly that the angel had not only the gospel, but that he had preached the gospel. Many indeed have the gospel, but it is silent and written in books. The gospel must be shown forth and proclaimed. He declares also to whom the gospel must be uttered and preached: to the inhabitants of the earth; for it must be cried out to those who are drowned in earthly matters, and they must be raised out of their sleep. And after his manner and imitation of blessed Daniel in the seventh chapter, he lists nations, kindreds, tongues, and peoples, thus signifying that the gospel shall be preached throughout the whole world. Which thing the Lord said also should come to pass in Matthew 24, and then that the end should come. And we see at this day that the gospel has in a manner thundered throughout the whole world. And here I give warning lest any deceive himself. The Apostle in 1 Timothy 3 and 1 Colossians says that the gospel was preached throughout the whole world in his time. However, not all men had then received it, but a few. Do not therefore imagine with yourself that the gospel is not preached unless all receive it. They are deceived who promise themselves before the judgment a concord of all nations, for it is written that there should be one shepherd and one flock. For the same was accomplished when the Lord prepared to himself one church from the Jewish synagogue and the dispersion of the gentiles, of which Christ is head and pastor, and Antichrist shall at the length by his last coming be abolished. Therefore he shall always resist Christ.

Furthermore, where he sees and hears this Angel preach the gospel with a loud voice, he means that the preachers shall with great constancy and frankness, also with shrill voices and utmost earnestness preach the gospel against Antichrist. And we see at this day that the more cruelly the faithful are grieved and oppressed, the more fervently and louder they cry, and that they also are called clamorous criers.

\index{Papacy}Moreover, he includes in a brief summary what things are to be presented in the preaching of the gospel, especially in the last times. First, he says: fear God. The fear of God is the beginning of wisdom; therefore, not fearing God is the beginning of foolishness and of all errors. The fear of God has nothing in common with the fear of the world. The godly man is not afraid of God as a guilty servant fears his master, and that punishment more than his master, whom he hates rather. For the fear of God has the reverence and love of God. It attributes to God supreme majesty, embraces faith, and has a faithful care, whereby it waits upon God, worships, praises, and professes him. Doubtless, because we fear men more than God, we fear more the Pope and the malice and hatred of him and his, therefore we do not execute justice uprightly, nor profess the truth frankly, neither yet set forth the gospel. But the Lord in the Gospel says: fear them not, which may kill the body and have no power over the soul; rather fear him, which condemns both body and soul to hellfire. Certainly, the fear of God is not only the beginning but also the bond of all virtue. Hereafter we shall hear that the fearful shall be cast into hell, with the beast and with the false prophet. Therefore let God be our fear, as Isaiah teaches in chapter 8. Let us fear God for our sins committed. Here few are afraid; but many are afraid to speak the truth, to maintain godliness, and to rebuke wickedness.

Secondly, the preaching of the Gospel comprehends the honor of God. For he says, “And give him honor.” And you do not separate the Son from the Father. For he, in John 5, says thus: “The Father has given all judgment to the Son, that all should honor the Son, as they honor the Father. He who does not honor the Son does not honor the Father who sent him.” And indeed, the Father cannot be honored but by the Son. For we honor him when we believe him to be true and receive Christ, the Son of God, as the only righteousness and perfection of all the faithful. By faith, therefore, we chiefly honor God, then revering him by faithful obedience and walking in his commandments. John, in his Canonical writings, says that he who does not believe in the Son makes God a liar—see how greatly one may dishonor God—who does not believe the testimony which God has testified about his Son. And this is the testimony, that God has given to us eternal life, and this life is in his Son. He who has the Son has life; he who does not have the Son does not have life. We are therefore forbidden, elsewhere than in Christ alone, to seek life and all goodness. But the papists honor the Pope, and his authority, his saints also, and honor God alone. They carve into their tombs, \textit{Soli deo gloria}—to God alone be glory—but yet in the meantime they persecute those who will not ascribe the glory due to God alone to their foolish trifles. But the Gospel will cry out that to God alone all glory is due.

Hereunto is added a spur, which may prick them to fear and glorify God: for the hour of his judgment is come. The Gospel therefore in the latter days shall impress upon men the final judgment. This has a wonderful effect in obtaining a reformation of life. And it is purposely said, it is come: and not, the hour of his judgment shall come. For so is the certainty of his judgment expressed, and we are warned to look for that same day every moment. The Apostle used the same argument in the 17th chapter of the Acts, addressed to the Athenians, and to the Corinthians in the 2nd Epistle, 5th chapter. Let us remember, O brethren, that strict judgment; let us amend our faith and manners, and all things that do not agree with the gospel. For certainly we shall die, certainly we shall be judged. But then, when we shall promise ourselves peace, sudden destruction shall come. Watch.

Finally, the Gospel teaches us to worship God alone. Do not the faithful then worship idols, wherever they may be erected? They do not worship the Pope, who is overwhelmed with wickedness; much less do they kiss, and by kissing worship his ungracious and foul feet. They do not worship the god Mauzim of the wafer makers, the god in the box, which is worshipped in palaces and churches as shut up in the pixe. They do not worship saints, but God alone. Therefore, lift up your hearts to heaven and worship. We have here on Earth wonderful works which may move us to worship this God alone. He is the maker of heaven and Earth, and of the sea. Who is greater? Who is mightier? Therefore, worship him, the true God, as in Matthew 4. He adds here fountains of water, for the miracle and benefit of water is great. For if we consider the original spring, substance, pleasantness, and usefulness of fountains, we shall be compelled to wonder. God be praised.
